\noindent\textbf{Motivation.} This book develops a compositional view of digital systems in which structure is not an afterthought but a source of mathematical leverage. The central themes are compositionality, closure, and symmetry: how complex systems are built from smaller ones, which families of operations remain stable under composition, and which invariances reveal hidden structure. These ideas recur across Boolean functions, circuits, categorical semantics, and related computational models, and they provide a common language for reasoning about hardware, software, and other discrete systems.

\medskip
\noindent\textbf{Audience and prerequisites.} The intended reader is comfortable with discrete mathematics, basic category theory, and digital logic. Familiarity with sets, functions, relations, Boolean algebra, and elementary proof techniques will be useful throughout. Some chapters also assume a working knowledge of algebraic structures and diagrammatic reasoning, but the exposition is designed to build these ideas incrementally rather than presuppose advanced specialization.

\medskip
\noindent\textbf{How to use this book.} The book is organized around a proof-as-builds workflow: mathematical claims are treated as constructions that can be checked, refined, and reused. Correspondingly, code is treated as witness material for the claims it implements, and each chapter is intended to leave behind concrete artifacts---definitions, derivations, examples, and verification traces---that can be inspected independently. Readers may proceed linearly, or use the chapter structure as a reference map: the early chapters establish the foundational language, the middle chapters develop modular constructions and feedback, and the later chapters apply the same principles to symmetry analysis, robustness, programs, dataflow, quantum-flavored models, cryptography, learning, and algorithms.

\medskip
\noindent The guiding expectation is that each result should be understandable at multiple resolutions: as an abstract statement, as a diagrammatic or algebraic derivation, and as an artifact that can be reproduced or adapted. If you are interested primarily in applications, the later chapters can be read with selective consultation of the foundations; if you are interested in the theory, the book is designed to reward a careful sequential reading.

\medskip
\noindent In short, this book aims to provide a unified framework for understanding how closure and symmetry organize digital composition, together with a disciplined method for turning that understanding into reusable technical artifacts.
